\section{Реализованные алгоритмы}
\label{sec:Chapter4} \index{Chapter4}
 
% ==============================
% 5. Реализованные алгоритмы
% ==============================

В данной главе описываются три основных механизма, реализованных в рамках курсовой работы: 
горизонтальное разрезание таблиц, многопроходный алгоритм \texttt{FFD} (First Fit Decreasing) и алгоритм 
\texttt{FFLS} (First Fit by Level and Size).  
Первый является вспомогательным и применяется, когда таблица не помещается целиком в одну стадию. 
Второй представляет собой авторскую модификацию классического \texttt{FFD}, позволяющую избегать тупиков 
за счёт многопроходного цикла. Третий — известная из литературы эвристика, 
реализованная для сравнения. Все алгоритмы учитывают аппаратные ограничения, а также типы зависимостей между таблицами .

\subsection{Вычисление потребности в блоках памяти}

Перед размещением для каждой логической таблицы $t$ необходимо определить, сколько физических блоков 
памяти (SRAM или TCAM) потребуется для хранения всех её записей.

Обозначим:
\begin{itemize}
	\item $w(t)$ — ширина ключа таблицы (бит);
	\item $e(t)$ — максимальное количество записей в таблице;
	\item $W_{\text{mem}}$ — ширина одного блока памяти (бит);
	\item $D_{\text{mem}}$ — глубина блока (количество ячеек, или записей, которое может хранить один блок).
\end{itemize}

Для точного поиска (exact) используется SRAM, для тернарного и префиксного (ternary, lpm) — TCAM.  
Количество блоков, необходимых для таблицы $t$, вычисляется по формуле:

\[
\text{blocks}(t) = \left\lceil \frac{w(t)}{W_{\text{mem}}} \right\rceil \;\times\; \left\lceil \frac{e(t)}{D_{\text{mem}}} \right\rceil.
\]

 где $\lceil w(t) / W_{\text{mem}} \rceil$ —  определяет, сколько блоков нужно для хранения одной записи (если ключ шире блока, он распределяется по нескольким блокам). Второй множитель — сколько таких «полосок» требуется для размещения всех записей. Полученное значение $\text{blocks}(t)$ используется во всех последующих алгоритмах как «размер» таблицы.

\subsection{Горизонтальное разрезание таблиц}

Горизонтальное разрезание применяется, когда $\text{blocks}(t)$ превышает максимальное количество блоков $C$, которое может быть выделено в одной стадии для данного типа памяти. Таблица разбивается на $k = \lceil \text{blocks}(t) / C \rceil$ фрагментов. Размер первого фрагмента равен $C$, последнего — $\text{blocks}(t) - (k-1)C$, все промежуточные — также $C$. Фрагменты размещаются в последовательных стадиях конвейера, начиная с некоторой начальной стадии $s$, которая определяется динамически в процессе работы алгоритмов размещения.

\noindent\textbf{Условия корректности разрезания:}
\begin{enumerate}
	\item Фрагменты одной таблицы должны занимать стадии с номерами $s, s+1, \dots, s+k-1$ без пропусков.  
	Это требование обусловлено аппаратной реализацией RMT: после неудачного поиска в стадии $i$ пакет переходит в стадию $i+1$, а не в произвольную.
	\item Если таблица $A$ является предком для таблицы $B$ по match- или reverse\_match-зависимости и хотя бы одна из них разрезана, то должно выполняться  
	\[
	\max_{\text{фрагменты } A} \text{stage} \;<\; \min_{\text{фрагменты } B} \text{stage}.
	\]  
	Если разрезана только $B$ (потомок), то достаточно $\text{stage}(A) < \min \text{stage}(B)$; если разрезана только $A$ — $\max \text{stage}(A) < \text{stage}(B)$.
\end{enumerate}

\subsection{Многопроходный алгоритм FFD}

Классический \texttt{FFD} сортирует таблицы по убыванию $\text{blocks}(t)$ и однократно размещает их в первую подходящую стадию. Его недостаток — возможные тупики при наличии match-зависимостей: если большая таблица-потомок обрабатывается раньше маленького предка, то предок впоследствии не может быть размещён.  

Предлагаемый многопроходный \texttt{FFD} устраняет эту проблему за счёт повторных проходов по ещё не размещённым таблицам.

На каждом проходе перебираются ещё не размещённые таблицы в порядке убывания размера. Для таблицы, у которой все предки уже размещены, запускается процедура размещения. Она пытается поместить первый фрагмент в существующие стадии, начиная с первой. Если подходящей стадии нет, создаётся новая. После размещения фрагмента остаток таблицы,если он есть, помещается в следующую по номеру стадию.

\begin{algorithm}[H]
	\caption{Многопроходный FFD с разрезанием}
	\begin{algorithmic}[1]
		Input: Множество таблиц $T$, ёмкость стадии $C$, зависимости \\
		Output: Размещение $F$ (фрагменты по стадиям) или ошибка
		\STATE вычислить $\text{blocks}(t)$ для всех $t \in T$
		\STATE отсортировать $T$ по убыванию $\text{blocks}(t)$
		\STATE $S \gets \emptyset$ (список стадий), $\text{placed} \gets \emptyset$
		\WHILE{$T \setminus \text{placed} \neq \emptyset$}
		\STATE $\text{progress} \gets \text{False}$
		\FOR{each $t \in T \setminus \text{placed}$ в порядке сортировки}
		\IF{$\text{pred}(t) \subseteq \text{placed}$} 
		\STATE попытаться разместить все фрагменты $t$ в существующих или новых стадиях (First Fit с последовательными номерами)
		\IF{размещение удалось}
		\STATE добавить фрагменты в $F$, $\text{placed} \gets \text{placed} \cup \{t\}$, $\text{progress} \gets \text{True}$
		\ENDIF
		\ENDIF
		\ENDFOR
		\IF{$\neg \text{progress}$} \RETURN ошибка (тупик) \ENDIF
		\ENDWHILE
		\RETURN $F$
	\end{algorithmic}
\end{algorithm}




\subsection{Алгоритм FFLS (First Fit by Level and Size)}

Алгоритм \texttt{FFLS}~\cite{jose2015compiling} является однопроходным и основан на топологическом упорядочении таблиц. Он гарантирует, что ни одна таблица не будет обработана раньше своих предков, что исключает тупики, связанные с порядком размещения (однако из-за жадной стратегии возможны проблемы с упаковкой памяти, особенно при разрезании).

\subsubsection{Вычисление уровней}

Для каждой таблицы $t$ вычисляется уровень $\text{level}(t)$ как длина самого длинного пути в графе зависимостей от любого источника (таблицы без предков). Формально:
\[
\text{level}(t) = 
\begin{cases}
	0, & \text{если } \text{pred}(t) = \varnothing;\\
	\max\limits_{p \in \text{pred}(t)} (\text{level}(p) + 1), & \text{иначе}.
\end{cases}
\]
Для любого ребра зависимости $p \to t$ выполняется $\text{level}(p) < \text{level}(t)$.


\begin{algorithm}[H]
	\caption{FFLS с разрезанием}
	\begin{algorithmic}[1]
		\REQUIRE Множество таблиц $T$, ёмкость стадии $C$, зависимости
		\ENSURE Размещение $F$
		\STATE вычислить $\text{blocks}(t)$ и $\text{level}(t)$ для всех $t$
		\STATE отсортировать $T$ по $(\text{level}(t), -\text{blocks}(t))$ (возрастание уровня, убывание размера)
		\STATE $S \gets \emptyset$ (список стадий)
		\FOR{each $t \in T$ в отсортированном порядке}
		\STATE разместить фрагменты $t$ методом First Fit (аналогично шагу размещения в многопроходном FFD, но без внешнего цикла по проходам)
		\IF{размещение невозможно} \RETURN ошибка \ENDIF
		\ENDFOR
		\RETURN $F$
	\end{algorithmic}
\end{algorithm}


Сначала вычисляются уровни таблиц, затем выполняется сортировка: сначала по возрастанию уровня, потом по убыванию $\text{blocks}(t)$. Это гарантирует, что все предки встречаются раньше потомков. Далее таблицы обрабатываются однократно в этом порядке. Для каждой таблицы применяется та же процедура размещения с поддержкой разрезания и проверкой зависимостей, что и в многопроходном FFD. Благодаря корректному порядку нет необходимости в повторных проходах. 






